\documentclass[10pt, letterpaper]{article}

% Packages:
\usepackage[
    ignoreheadfoot, % set margins without considering header and footer
    top=1 cm, % seperation between body and page edge from the top
    bottom=1 cm, % seperation between body and page edge from the bottom
    left=2 cm, % seperation between body and page edge from the left
    right=2 cm, % seperation between body and page edge from the right
    footskip=1.0 cm, % seperation between body and footer
    % showframe % for debugging 
]{geometry} % for adjusting page geometry
\usepackage{titlesec} % for customizing section titles
\usepackage{tabularx} % for making tables with fixed width columns
\usepackage{array} % tabularx requires this
\usepackage[dvipsnames]{xcolor} % for coloring text
\definecolor{primaryColor}{RGB}{0, 79, 144} % define primary color
\usepackage{enumitem} % for customizing lists
\usepackage{fontawesome5} % for using icons
\usepackage{amsmath} % for math
\usepackage[
    pdftitle={Rabin Dhamala's CV},
    pdfauthor={Rabin Dhamala},
    pdfcreator={Rabin Dhamala},
    colorlinks=true,
    urlcolor=primaryColor
]{hyperref} % for links, metadata and bookmarks
\usepackage[pscoord]{eso-pic} % for floating text on the page
\usepackage{calc} % for calculating lengths
\usepackage{bookmark} % for bookmarks
\usepackage{lastpage} % for getting the total number of pages
\usepackage{changepage} % for one column entries (adjustwidth environment)
\usepackage{paracol} % for two and three column entries
\usepackage{ifthen} % for conditional statements
\usepackage{needspace} % for avoiding page brake right after the section title
\usepackage{iftex} % check if engine is pdflatex, xetex or luatex

% Ensure that generate pdf is machine readable/ATS parsable:
\ifPDFTeX
    \input{glyphtounicode}
    \pdfgentounicode=1
    % \usepackage[T1]{fontenc} % this breaks sb2nov
    \usepackage[utf8]{inputenc}
    \usepackage{lmodern}
\fi



% Some settings:
\AtBeginEnvironment{adjustwidth}{\partopsep0pt} % remove space before adjustwidth environment
\pagestyle{empty} % no header or footer
\setcounter{secnumdepth}{0} % no section numbering
\setlength{\parindent}{0pt} % no indentation
\setlength{\topskip}{0pt} % no top skip
\setlength{\columnsep}{0cm} % set column seperation
\makeatletter

\titleformat{\section}{\needspace{4\baselineskip}\bfseries\large}{}{0pt}{}[\vspace{1pt}\titlerule]

\titlespacing{\section}{
    % left space:
    -1pt
}{
    % top space:
    0.3 cm
}{
    % bottom space:
    0.2 cm
} % section title spacing

\renewcommand\labelitemi{$\circ$} % custom bullet points
\newenvironment{highlights}{
    \begin{itemize}[
        topsep=0.10 cm,
        parsep=0.10 cm,
        partopsep=0pt,
        itemsep=0pt,
        leftmargin=0.4 cm + 10pt
    ]
}{
    \end{itemize}
} % new environment for highlights

\newenvironment{highlightsforbulletentries}{
    \begin{itemize}[
        topsep=0.10 cm,
        parsep=0.10 cm,
        partopsep=0pt,
        itemsep=0pt,
        leftmargin=10pt
    ]
}{
    \end{itemize}
} % new environment for highlights for bullet entries


\newenvironment{onecolentry}{
    \begin{adjustwidth}{
        0.2 cm + 0.00001 cm
    }{
        0.2 cm + 0.00001 cm
    }
}{
    \end{adjustwidth}
} % new environment for one column entries

\newenvironment{twocolentry}[2][]{
    \onecolentry
    \def\secondColumn{#2}
    \setcolumnwidth{\fill, 4.5 cm}
    \begin{paracol}{2}
}{
    \switchcolumn \raggedleft \secondColumn
    \end{paracol}
    \endonecolentry
} % new environment for two column entries

\newenvironment{header}{
    \setlength{\topsep}{0pt}\par\kern\topsep\centering\linespread{1.5}
}{
    \par\kern\topsep
} % new environment for the header

\newcommand{\placelastupdatedtext}{% \placetextbox{<horizontal pos>}{<vertical pos>}{<stuff>}
  \AddToShipoutPictureFG*{% Add <stuff> to current page foreground
    \put(
        \LenToUnit{\paperwidth-2 cm-0.2 cm+0.05cm},
        \LenToUnit{\paperheight-1.0 cm}
    ){\vtop{{\null}\makebox[0pt][c]{
        \small\color{gray}\textit{Last updated in September 2024}\hspace{\widthof{Last updated in September 2024}}
    }}}%
  }%
}%

% save the original href command in a new command:
\let\hrefWithoutArrow\href

% new command for external links:
\renewcommand{\href}[2]{\hrefWithoutArrow{#1}{\ifthenelse{\equal{#2}{}}{ }{#2 }\raisebox{.15ex}{\footnotesize \faExternalLink*}}}


\begin{document}
    \newcommand{\AND}{\unskip
        \cleaders\copy\ANDbox\hskip\wd\ANDbox
        \ignorespaces
    }
    \newsavebox\ANDbox
    \sbox\ANDbox{}

    % \placelastupdatedtext
    \begin{header}
        \textbf{\fontsize{24 pt}{24 pt}\selectfont Rabin Dhamala}

        \vspace{0.3 cm}

        \normalsize
        \mbox{{\color{black}\footnotesize\faMapMarker*}\hspace*{0.13cm}Basundhara, Kathmandu}%
        \kern 0.25 cm%
        \AND%
        \kern 0.25 cm%
        \mbox{\hrefWithoutArrow{rabinkmc@gmail.com}{\color{black}{\footnotesize\faEnvelope[regular]}\hspace*{0.13cm}rabinkmc@gmail.com}}%
        \kern 0.25 cm%
        \AND%
        \kern 0.25 cm%
        \mbox{\hrefWithoutArrow{tel:+977-9823189006}{\color{black}{\footnotesize\faPhone*}\hspace*{0.13cm}+977 9823189006}}%
        \kern 0.25 cm%
        \AND%
        \kern 0.25 cm%
        \mbox{\hrefWithoutArrow{https://rabinkmc.github.io/}{\color{black}{\footnotesize\faLink}\hspace*{0.13cm}rabinkmc.github.io}}%
        \kern 0.25 cm%
        \AND%
        \kern 0.25 cm%
        % \mbox{\hrefWithoutArrow{https://linkedin.com/in/rabin-dhamala-785581134/}{\color{black}{\footnotesize\faLinkedinIn}\hspace*{0.13cm}rabin-dhamala-785581134}}%
        % \kern 0.25 cm
        % \AND%
        % \kern 0.25 cm%
        % \mbox{\hrefWithoutArrow{https://github.com/yourusername}{\color{black}{\footnotesize\faGithub}\hspace*{0.13cm}yourusername}}%
    \end{header}

    \vspace{0.3 cm - 0.3 cm}

    \section{Experience}

        \begin{twocolentry}{
            \textit{Kathmandu, Nepal}    

            \textit{April 2024 - Sep 2024}
        }
            \textbf{Software Engineer}

            \textit{Uptechsys Pvt. Ltd}
        \end{twocolentry}
        \vspace{0.1 cm}
        \begin{onecolentry}
            \begin{highlights}
            \item Here, I primarily worked as a Django Backend Developer. Here, I deployed and maintained several projects. Notably, I worked on Hamro Academy, E-Nagari, Municipal Planning Management System. I also setup the blog infrastructure system. Additionally, I improved the integration and deployment workflow by integrating a testing environment and automating deployment for several of these projects. 
        \end{highlights}
        \end{onecolentry}
        \vspace{0.1 cm}
        \begin{twocolentry}{
        \textit{Remote, US}    
            
        \textit{Feb 2023 – Nov 2023}}
            \textbf{Software Engineer}
            
            \textit{Cleavr}
        \end{twocolentry}

        \vspace{0.10 cm}
        \begin{onecolentry}
            \begin{highlights}
            \item The work was entirely remote, requiring clear communication to ensure alignment on my plans, the team's priorities, and what was best for the company. The project involved a small, agile team of 2-3 members, where we handled everything from frontend and backend development to writing blogs and providing customer support. The domain required a solid understanding of DevOps practices, and over time, I gained valuable knowledge in areas such as \textbf{SSL}, \textbf{Nginx}, \textbf{DNS}, \textbf{monitoring tools}, and \textbf{system services}. Bash scripting was essential for setting up these services effectively, and the experience significantly deepened my technical expertise.

            \item The tech stack involved \textbf{AdonisJs} on the backend and \textbf{Vue} on the frontend following a monolithic architecture connected via \textbf{InertiaJs}. I worked on a variety of issues, such as providing support for new versions of database such as Postgres, Mariadb, fixing UI respnsiveness issues in frontend, integrating a complete \textbf{email change system}, enabling \textbf{server restart} via web portal, and adding feature such as \textbf{disk cleanup} etc. 
            \end{highlights}
        \end{onecolentry}


        \vspace{0.2 cm}

        \begin{twocolentry}{
        \textit{Kathmandu, Nepal}    
            
        \textit{Oct 2021 – March 2023}}
            \textbf{Software Engineer}
            
            \textit{Aayulogic Pvt.Ltd}
        \end{twocolentry}

        \vspace{0.10 cm}
        \begin{onecolentry}
            \begin{highlights}
            \item Here, I specialized in \textbf{Python}, working mainly in \textbf{Django} and \textbf{DRF} to build \textbf{APIs} for the frontend. I worked on attendance penalty system, optimizing payroll, transforming static recruitment process to dynamic. Technically, I gained knowledge in implementing and maintaining authentication system, \textbf{RBAC} system, learnt and used \textbf{Docker} in our project workflow, implemented caching, used multiprocessing to parallelize CPU intensive tasks, used threading (\textbf{Django Q}) for asynchronous tasks like email and notification. 
            \end{highlights}
        \end{onecolentry}
    
    \section{Education}
        \begin{twocolentry}{
            
            
        \textit{Nov 2016 –  Sept 2021}}
            \textbf{Pulchowk Campus, Tribhuvan University, Nepal}

            \textit{Bachelors in Electrical Engineering}
        \end{twocolentry}

        \vspace{0.10 cm}
        \begin{onecolentry}
            \begin{highlights}
                \item scored 72.5 \% overall
                \item Notable courseworks on C, C++, Numerical Methods, Microprocess 8085, 8086, System Reliability and Artificial Neural Network
            \end{highlights}
        \end{onecolentry}

    \section{Technical Growth}
            \begin{highlights}
            \item \textbf{Data structure and Algorithm}: After taking the break from work, I primarily focused on DSA, solving over 800 LeetCode problems on DSA. This is a self repository, \href{https://github.com/rabinkmc/dsa}{DSA}. The objective was to improve problem solving ability by working on DP problems, binary search algorithm, graph problems, and searching algorithms.
            \item \textbf{Codecrafters challenge}: I completed several codecrafters challenge, including building \href{https://github.com/rabinkmc/http-server-python}{http server}, \href{https://github.com/rabinkmc/regex-engine}{regex engine}, \href{https://github.com/rabinkmc/codecrafters-redis-python}{Redis}, \href{https://github.com/rabinkmc/dns-server-python}{DNS server}, \href{https://github.com/rabinkmc/docker-python}{Docker program to pull image from Docker hub}, \href{https://github.com/rabinkmc/git-python}{building your own Git} on early 2024. The main aim was to learn the underlying technology by implementing some part of the tool itself. I have spent the the last 3-4 months (Oct 2025 to Feb 2026 ) again working on \href{https://github.com/rabinkmc/redis-go}{Redis}, \href{https://github.com/rabinkmc/interpreters}{Interpreter}, creating your own \href{https://github.com/rabinkmc/shell-go}{Shell},. The aim was to be fairly comfortable with low level programming in "Go". 
            \item \textbf{Deep Learning Specialization course by Andrew Ng on Coursera}. I learnt about neural networks (CNN, RNN) and implemented them using Tensorflow. \href{https://coursera.org/share/32d36a0abf73bc14c097a05e7d1296fa}{Deep learning Specialization}
            \item \textbf{Career break since Oct 2024 }
                Took a planned year-long break to prepare for master’s studies. During this period, I deepened my understanding of operating systems, distributed systems, computer networks, and databases through self-study. Although I wasn’t able to pursue studies abroad as planned, the experience strengthened my technical foundation and reinforced my commitment to returning to professional software development with renewed focus and drive. However, it would be in my bad faith to paint a rosy picture of this but I believe I am a better SWE having gone through this period.
            \end{highlights}

    \section{Programming Skills}
    \begin{highlights}
    \item \textbf{Languages}: Python, Javascript, NodeJs, Golang, (C/C++, Bash scripting - basic)
    \item \textbf{Framework and Technologies}: Docker, Django, Vue, Redis, Nginx, Postgresql
    \end{highlights}

    % \section{AI and Development Tools}
    %     \begin{highlights}
    %         \item Use AI-assisted tools such as ChatGPT, GitHub Copilot, and Claude Code to accelerate development and debugging.
    %         \item Emphasize strong computer science fundamentals to ensure correctness, originality, and deep understanding beyond AI-generated code.
    %         \item View AI tools as productivity aids rather than replacements for critical thinking, and apply them with caution and intent.
    %     \end{highlights}
\end{document}
